\section{Ancestro Comun mas Bajo (LCA)}

Es un algoritmo que te permite encontrar el ancestro comun mas bajo de dos nodos en un arbol. Por conveniencia, suele tomar como raiz, el primer nodo (i=0). Se suele utilizar para conocer la distancia entre dos nodos en un arbol usando la formula de $D(x)+D(y)-D(lca(x,y))$ asumiendo $D$ como la profundidad del nodo (distancia desde la raiz).

Su algoritmo es uno de los mas complejos y dificiles de entender. Aca se dira una forma sumamente simplificada de como (se cree que) funciona.
\begin{itemize}
    \item Se inicializa una matriz $F[k][2^k]$ donde $F[x][n]$ es el padre $2^x$ del nodo $n$.
    \item Se inicializa un Array $D[2^k]$ con la profundidad de cada nodo.
    \item Se inicializa la profundidad del nodo 0 en 0 (cero) y su nodo padre $F[0][0]=-1$ marcando que es el raiz.
    \item Se recorre el arbol, empezando desde el nodo 0, guardando para cada nodo, su padre y su profundidad.
    \item Por cada potencia de dos y por cada nodo, verifico si el padre en la potencia $k-1$ no existe, entonces el padre en la potencia $k$ tampoco existe. Si, en cambio, el padre la potencia $k-1$ existe, el padre en la potencia $k$ es el padre en la potencia $k-1$ del padre en la potencia $k-1$ de $x$.
\end{itemize}
Toda esta informacion, se precalcula y se almacena, los siguiente pasos son para encontrar el LCA de dos nodos.
\begin{itemize}
    \item Por cada potencia de dos, empezando desde la mas alta posible, verifico si el nodo con mayor profundidad puede subir esas posiciones y quedar en una profundidad igual  o mayor a la del otro nodo. Hasta que se igualen las profundidades.
    \item Se verifica si el nodo padre de que era menor, es el otro nodo, de ser asi, devolvemos eso, sino continuamos.
    \item Por cada potencia de dos, empezando desde la mas alta posible, verifico si el padre, en dicha potencia, de ambos, es distinto, entonces subimos ambos nodos hasta dichos padres.
    \item Finalmente, tendremos los dos primeros ancestros distintos, por lo que devolvemos su padre.
\end{itemize}

\begin{lstlisting} [style=cpp]
const int K = 10;
vector<int> g[1<<K]; int N; // K de tal manera que 2^K>=N
int F[K][1<<K], D[1<<K]; // F[i][j] es el padre 2^i del nodo j
// F[i][j] donde 'i' representa la potencia de dos en la que 
//se encuentra el padre de 'j'

void lca_dfs(int x){
    for(int hijo : g[x]){
        // Si es el padre, se continua
        if(hijo==F[0][x]) continue;
         // Se marca x como el padre del hijo
        F[0][hijo]=x;
        // La distancia del hijo es la del padre mas 1
        D[hijo]=D[x]+1;
        // Se recorren los hijos del nodo hijo
        lca_dfs(hijo); 
    }
}

void lca_init(){
    D[0]=0;F[0][0]=-1;
    lca_dfs(0); // Se recorre el arbol
    // Para cada potencia de 2 y para cada nodo
    fore(k,1,K) fore(x,0,N){ 
        // Si el padre que esta en la potencia 'k-1' no existe, 
        // entonces el padre en la potencia k, tampoco existe
        if(F[k-1][x]<0) F[k][x]=-1; 
        // Si no, el padre es el padre 'k-1' del padre 'k-1' de x
        else F[k][x] = F[k-1][F[k-1][x]]; 
    }
}

int lca(int x, int y){
    // Nos aseguramos de que el x sea el de mayor profundidad
    if (D[x]<D[y]) swap(x,y); 
    // Buscamos el padre del nodo x que esta en la misma 
    // profundidad que y.
    for(int k=K-1;k>=0;--k) if(D[x]-(1<<k)>=D[y]) x=F[k][x]; 
    if (x==y) return x; // Si son el mismo nodo, devolvemos x
    // Buscamos los primeros padres distintos de ambos nodos
    for(int k=K-1;k>=0;--k) if(F[k][x]!=F[k][y]) x=F[k][x], y=F[k][y];
    return F[0][x]; // Devolvemos su padre.
}
\end{lstlisting}


\section{LCA con binary lifting}
\whenuse{Encontrar el ancestro com\'un m\'as cercano, calcular distancias o subir
$k$ niveles en un \'arbol enraizado.}

El LCA de dos v\'ertices es el punto donde se unen sus caminos hacia la
ra\'iz. Se precomputa \texttt{up[k][v]}, el ancestro ubicado $2^k$ niveles
arriba de $v$, para realizar saltos de $1,2,4,8,\ldots$ niveles.

Para un \'arbol sin pesos:
\[
dist(a,b)=depth[a]+depth[b]-2\,depth[LCA(a,b)].
\]

\complexity{Preproceso $O(n\log n)$; consulta y salto $O(\log n)$;
memoria $O(n\log n)$.}

\begin{lstlisting}[style=cpp]
struct LCA {
    int n, levels;
    vector<int> depth;
    vector<vector<int>> up; // up[k][v]: ancestro 2^k de v
    LCA(const vector<vector<int>> &tree, int root = 0)
        : n(tree.size()) {
        assert(n > 0 && 0 <= root && root < n);
        levels = 1;
        while ((1 << levels) <= n) ++levels;
        depth.assign(n, 0);
        vector<int> parent(n, -1);
        // Enraizar el arbol y calcular padres y profundidades.
        queue<int> q;
        parent[root] = root;
        q.push(root);
        while (!q.empty()) {
            int v = q.front();
            q.pop();

            for (int to : tree[v]) {
                if (parent[to] != -1) continue;
                parent[to] = v;
                depth[to] = depth[v] + 1;
                q.push(to);
            }
        }
        // Construir la tabla de saltos.
        up.assign(levels, vector<int>(n));
        up[0] = parent;
        for (int k = 1; k < levels; ++k)
            for (int v = 0; v < n; ++v)
                up[k][v] = up[k - 1][up[k - 1][v]];
    }
    // Sube distance niveles desde v.
    int jump(int v, int distance) const {
        for (int k = 0; k < levels; ++k)
            if (distance & (1 << k))
                v = up[k][v];
        return v;
    }
    // Ancestro comun mas cercano de a y b.
    int query(int a, int b) const {
        // Igualar profundidades.
        if (depth[a] < depth[b]) swap(a, b);
        a = jump(a, depth[a] - depth[b]);
        if (a == b) return a;
        // Subir ambos sin alcanzar todavia el LCA.
        for (int k = levels - 1; k >= 0; --k) {
            if (up[k][a] != up[k][b]) {
                a = up[k][a];
                b = up[k][b];
            }
        }
        return up[0][a];
    }
    // Cantidad de aristas en el camino entre a y b.
    int distance(int a, int b) const {
        int ancestor = query(a, b);
        return depth[a] + depth[b] - 2 * depth[ancestor];
    }
};
vector<vector<int>> tree(n);
auto add_edge = [&](int u, int v) {
    tree[u].push_back(v);
    tree[v].push_back(u);
};
add_edge(0, 1);
add_edge(0, 2);
add_edge(1, 3);
add_edge(1, 4);
LCA lca(tree, 0);
cout << lca.query(3, 4) << '\n';    // 1
cout << lca.jump(3, 2) << '\n';     // 0
cout << lca.distance(3, 2) << '\n'; // 3
\end{lstlisting}


\warning{El grafo debe ser un \'arbol conectado. \texttt{jump(v,k)} supone
$0\le k\le depth[v]$; la distancia devuelta cuenta aristas.}